\section{Introduction}
\label{sec:intro}

The training pipeline of modern foundation models is traditionally partitioned into distinct sequential phases: self-supervised Next-Token Prediction (NTP) pre-training over large text corpora, followed by Supervised Fine-Tuning (SFT), and ending with Reinforcement Learning (RL) alignment \citep{bansal2026rl, huang2025reinforcement}. While this pipeline has yielded remarkable advances, applying RL strictly as a post-training stage introduces fundamental bottlenecks: late-stage RL primarily serves to "sharpen" probability distributions around high-reward modes, potentially compressing policy entropy and constraining language plasticity.

Recent studies demonstrate that introducing RL early or at intermediate pre-training checkpoints—termed "RL excursions"—can broaden output distribution coverage and foster diverse solution strategies \citep{bansal2026rl}. However, existing methods employ static or heuristic intervention timing, applying RL at pre-determined step intervals irrespective of whether the model checkpoint currently possesses the internal representations needed to profit from policy optimization.

In this work, we formulate the dynamic scheduling of reinforcement learning as a continuous feedback-control problem driven by checkpoint readiness signals. We ask four primary research questions:
\begin{enumerate}
    \item Which measurable properties of intermediate language-model checkpoints predict subsequent reinforcement-learning improvement?
    \item Can those properties distinguish between capability expansion, sharpening, stagnation, and interference?
    \item How can an adaptive training scheduler exploit these signals to allocate compute dynamically among NTP, SFT, and RL?
    \item Does adaptive scheduling improve compute-normalized performance compared to fixed baseline schedules?
\end{enumerate}

To answer these questions, we perform a systematic empirical checkpoint plasticity study and introduce the \textbf{Capability-Aware Reinforcement Learning Scheduler (CARLS)}. CARLS evaluates real-time readiness indicators—most notably the gradient alignment between NTP pre-training loss and reward-driven policy gradients $\cos(\mathbf{g}_{\text{NTP}}, \mathbf{g}_{\text{RL}})$, policy entropy, and baseline pass@k accuracy—to dynamically adjust resource allocation among training objectives.
